\documentclass[10pt,a4paper]{article}
\usepackage[utf8]{inputenc}
\usepackage{ctex} % 支持中文排版
\usepackage[left=18mm, right=18mm, top=18mm, bottom=18mm]{geometry}
\usepackage{titlesec}
\usepackage{enumitem}
\usepackage{amssymb}
\usepackage{color}
\usepackage{hyperref}

% 颜色定义 - 专业深蓝色调（适合电子信息/物理等理工科专业）
\definecolor{primarycolor}{RGB}{30, 65, 110}
\definecolor{graycolor}{RGB}{100, 100, 100}

% 页面样式调整
\pagestyle{empty}
\linespread{1.2}

% 标题格式设置
\titleformat{\section}{\large\bfseries\color{primarycolor}}{}{0em}{}[\titlerule]
\titlespacing{\section}{0pt}{10pt}{6pt}

\titleformat{\subsection}{\normalsize\bfseries}{}{0em}{}
\titlespacing{\subsection}{0pt}{4pt}{2pt}

\begin{document}

% ==================== 个人信息 ====================
\begin{center}
    {\LARGE \bfseries \color{primarycolor} [姓名] 个人简介} \\
    \vspace{4pt}
    \small
    姓名：[姓名] \quad |性别：男 \quad | \quad 年龄：20岁 \quad | \quad 籍贯：内蒙古包头市
    \vspace{2pt}
    | \quad 电话：13664770236 
\end{center}

% ==================== 教育背景 ====================
\section{教育背景}
\textbf{南开大学} 
\hfill \textbf{2024年8月入学} \\
物理学 | 物理科学学院 \hfill 大一 \\
电子信息科学与技术 | 电子信息与光学工程学院 \hfill 大二在读 \\
\begin{itemize}[leftmargin=15pt, topsep=2pt, itemsep=0pt]
    \item \textbf{高考成绩}：总分 661/750（物化生），理科综合 270 分，数学 134 分
    \item \textbf{英语特长}：高考英语 145 分；高考一模、二模、三模英语单科成绩分别位列\textbf{全市第 7、第 2、第 3 名}
    \item \textbf{大学英语}：2025年6月大学英语六级（CET-6）成绩 \textbf{582分}
    \item \textbf{荣誉职责}：高中物理课代表
\end{itemize}


% ==================== 教学与实践经历 ====================
\section{家教经历}
\textbf{高中英语 天津 线下} \hfill \textbf{2025年3月至2025年8月（我转校区后）} \\
\begin{itemize}[leftmargin=15pt, topsep=2pt, itemsep=0pt]
    \item \textbf{辅导对象}：初始英语成绩67/150，南开大学附属中学
    \item \textbf{辅导时间}：每周六上午2个小时
    \item \textbf{辅导成果}：英语成绩从60分不稳定提升至90分左右，由于我从八里台校区转至津南校区，且学业压力加大，经过与家长沟通停止了家教进程
\end{itemize}

\textbf{高中英语（线上）} \hfill \textbf{2024年 - 2025年高考} \\
\begin{itemize}[leftmargin=15pt, topsep=2pt, itemsep=0pt]
    \item \textbf{辅导对象}：同校高二升高三学弟（初始英语成绩处于年级垫底水平，约 60 分左右）
    \item \textbf{辅导策略}：运用个人独创的英语学习方法论，每次约 1 小时网课
    \item \textbf{辅导成果}：学生英语成绩从 70 分逆袭提升至 \textbf{100分左右}，2025高考英语106分
\end{itemize}

\textbf{初中高中 物理 } \hfill \textbf{2024年至今} \\

\begin{itemize}[leftmargin=15pt, topsep=2pt, itemsep=0pt]
    \item \textbf{辅导对象}：北京初三男生，基础较薄弱，初始物理成绩仅为 40-50 分水平
    \item \textbf{辅导策略}：初三开始，我针对他读题水平差，题意理解和公式理解不到位的缺点进行指导，加强练习
    \item \textbf{辅导成果}：初三期末物理成绩提升至 \textbf{83/100分}；升入高中后，物理成绩长期\textbf{稳定在 80（满分100）}水平
\end{itemize}


% ==================== 核心优势与方法论 ====================
\section{核心优势与个人技能}
\begin{itemize}[leftmargin=15pt, topsep=2pt, itemsep=2pt]
    \item \textbf{学科方法论}：沉淀并总结出一套自己的的\textbf{英语、物理、数学学科学习方法论}，自认为有对于学生问题的洞察力。、
    \item \textbf{专业技能}：英语水平很有自信，有大学物理学与电子信息科学的学习背景，扎实的高中物理，数学基础
\end{itemize}

% ==================== 竞赛与学术规划 ====================
\section{暑假时间安排}
\textbf{天津市 TI杯 全国大学生电子设计竞赛} \hfill \textbf{2026年暑假} \\
\begin{itemize}[leftmargin=15pt, topsep=2pt, itemsep=0pt]
    \item 2026年7月2日放假，暑假留校跟组科研，周末全天以及周中晚上时间可以自由安排。开学后周中可以网课，周末可以线下。
\end{itemize}

\end{document}